%! TeX program = lualatex
%! TEX options = -synctex=1 -interaction=nonstopmode -file-line-error --shell-escape "%DOC%"
\documentclass[10pt,a5paper]{mszalik}

\begin{document}

\thispagestyle{empty}

\begin{center}
	\vspace*{0.5cm}


	\bfseries\scshape
	\huge MSZA O NAJŚWIĘTSZEJ\\MARYI PANNIE

    \vspace*{0.5cm}

    \huge RORATE

    \vspace*{0.5cm}

	\vfill

	\begin{figure}[h]
		\centering
		
\includegraphics[width=0.45\linewidth]{logo.png}
	\end{figure}

	\vfill

    \vspace*{0.5cm}

    \bigskip

    {\large Duszpasterstwo Wiernych Tradycji Łacińskiej\\ w Archidiecezji Wrocławskiej}

\end{center}

\newpage

\vspace*{2cm}

\vspace*{\fill}

\hrulefill

{\footnotesize

    Śpiewy i teksty liturgiczne pochodzą z:

	\begin{itemize}[leftmargin=*]
		\item \textit{Liber Usualis} (1961),
        \item \textit{Mszału Rzymskiego} autorstwa benedyktynów tynieckich (Poznań 1963).
	\end{itemize}

	W niektórych przypadkach wykorzystano cyfrową wersję tekstów
	opublikowaną na portalach \textit{divinumofficium.com} i \textit{missalemeum.com}.

    \medskip

    Zapis gabc neumów: Andrew Hinkley et al. \textit{gregobase.selapa.net}.

	\bigskip

	Opracowanie: Mateusz Burszewski, Piotr Matuszak\\
	Skład i łamanie: Michał Siemaszko, Piotr Matuszak\\

    2025, wersja df\\

	\texttt{tradycja.archidiecezja.wroc.pl}
}


%%%%%%%%%%%%%%%%%%%%%%%%%%%%%%%%%%%%%%%%%%%%%%%%%%%%%%%%%%%%%%%%%%%%%%%%%%

\newpage

\fancyhead[CO]{MSZA O NAJŚWIĘTSZEJ MARYI PANNIE}
\fancyhead[CE]{RORATE}

% \begin{figure}[h]
%     \centering
%     \includegraphics[width=\linewidth]{Figures/in_conceptione_immaculata_bmv_missale_romanum_1962.png}
% \end{figure}

\section{MSZA ŚWIĘTA}

\oremus{Introit (Iz 45,8; Ps 84,2)}{
    Roráte, caeli, désuper, et nubes pluant justum: aperiátur terra, et gérminet Salvatórem\\\\

	Benedixísti, Domine, terram tuam: avertísti captivitátem Jacob. \medskip

	Glória Patri\dots \\

    \textit{Roráte, coeli\dots}
    }{
    Niebiosa, spuśćcie rosę z góry, a obłoki niech zleją z deszczem Sprawiedliwego.
    Niech się otworzy ziemia i zrodzi Zbawiciela.\\

    Łaskawy okazałeś się, Panie, dla Twej ziemi, na dobre odmieniłeś los Jakuba.\medskip

	Chwała Ojcu\dots\\

    \textit{Niebiosa\dots}
}

\bigskip

\grechangestaffsize{15}
\grecommentary[0.15cm]{Missa Cum iubilo (IX)}
\gregorioscore{Melody/ky--kyrie_ix--solesmes.1.gabc}

\grecommentary[0.15cm]{Missa De Angelis (VIII)}
\gregorioscore{Melody/ky--gloria_viii--solesmes.1.gabc}

\medskip

\oremus{Kolekta}{
    Deus, qui de beátæ Maríæ Vírginis útero Verbum tuum, Angelo nuntiánte, carnem suscípere voluísti:
    præsta supplícibus tuis; ut, qui vere eam Genetrícem Dei crédimus, eius apud te intercessiónibus adiuvémur.
    Per eúndem Dominum\dots
}{
    Boże, któryś zrządził, że za zwiastowaniem anielskim w żywocie Najświętszej Maryi Panny Słowo Twoje Ciałem się
    stało, pokornie Cię błagamy, spraw, aby wspomagało nas wstawiennictwo Tej, w której Boże Macierzyństwo
    wierzymy.
    Przez tegoż Pana\dots
}

\newpage

\proroctwo{Lekcja (Iz 7,10-15)}

Z księgi proroka Izajasza

W one dni: Rzekł Pan do Achaza, mówiąc: «Proś Pana Boga twego o znak dla siebie:
w głębinie piekła lub na wysokości w górze». I rzekł Achaz: «Nie będę prosił i nie będę kusił Pana».
I rzekł: «Słuchajcie więc, domie Dawida: Czyż to mało dla was dręczyć ludzi, że dręczycie także i Boga mojego?
Przeto sam Pan da wam znak: Oto Panna pocznie i porodzi Syna i nazwą imię Jego Emmanuel.
Masło i miód będzie pożywał, by umiał zło odrzucać, a wybierać dobro».

\oremus{Graduał (Ps 23,7.3-4)}{
    Tóllite portas, príncipes, vestras: et elevámini, portæ æternáles: et introíbit Rex glóriæ.

    Quis ascéndet in montem Dómini? aut quis stabit in loco sancto ejus? Innocens mánibus et mundo corde.
}{
    Bramy, podnieście swe szczyty i rozstąpcie się, prastare podwoje, aby mógł wkroczyć Król chwały.

    Kto wstąpi na górę Pana lub stanie na świętym Jego miejscu? Człowiek rąk nieskalanych i czystego serca.
}

\oremus{Alleluja (Łk 1,28)}{
    Allelúja, allelúja. Ave, María, grátia plena; Dóminus tecum: benedícta tu in muliéribus. Alleluia.
}{
    Alleluja, alleluja. Zdrowaś, Maryjo, łaski pełna, Pan z Tobą, błogosławionaś Ty między niewiastami. Alleluja.
}

% \vfill

% \centerline{\pgfornament[width=5cm]{68}}

% \vfill

% \newpage

\proroctwo{Ewangelia (Łk 1:26-38)}

Ciąg dalszy \ding{64} Ewangelii według Łukasza.

Onego czasu: Posłany jest Anioł Gabriel od Boga do miasta galilejskiego, które zwano Nazaret,
do Panny poślubionej mężowi, któremu było na imię Józef, z domu Dawidowego, a imię Panny Maryja.
I wszedłszy do Niej Anioł rzekł: «Bądź pozdrowiona, łaski pełna, Pan z Tobą, błogosławionaś Ty między niewiastami».
A Ona, gdy posłyszała, zatrwożyła się na słowa jego i rozważała, cóż by to było za pozdrowienie.
I rzekł Jej Anioł: «Nie lękaj się, Maryjo, albowiem znalazłaś łaskę u Boga!
Oto poczniesz w łonie i porodzisz syna, i nadasz Mu imię Jezus. Ten będzie wielkim i Synem Najwyższego nazwany będzie.
I da Mu Pan Bóg stolicę Dawida, ojca Jego, i będzie królował nad domem Jakubowym na wieki,
a królestwu Jego nie będzie końca».
A Maryja rzekła do Anioła: «Jakże się to stanie, skoro męża nie znam?»
I odpowiadając Anioł rzekł Jej: «Duch Święty zstąpi na Ciebie i moc Najwyższego Cię zacieni.
Przeto i co się z Ciebie narodzi Święte, Synem Bożym będzie nazwane.
A oto Elżbieta, krewna Twoja, i ona poczęła dziecię w starości swojej, a jest to już szósty miesiąc dla tej,
którą zowią niepłodną: bo nie masz nic niemożliwego u Boga». I rzekła Maryja:
«Oto ja służebnica Pańska, niechaj mi się stanie według słowa twego».


\newpage

% \grechangestaffsize{16}
% \grecommentary[0.15cm]{Credo (IV)}
% \gregorioscore{Melody/ky--credo_iv--solesmes.1.gabc}
% \grechangestaffsize{17}


\oremus{Ofertorium (Łk 1,28.42)}{
    Ave, María, grátia plena; Dóminus tecum: benedícta tu in muliéribus, allelúja.
}{
    Zdrowaś Maryjo, łaski pełna, Pan z Tobą, błogosławionaś Ty między niewiastami, alleluja.
}

\oremus{Sekreta}{
    In méntibus nostris, quǽsumus, Dómine, veræ fídei sacraménta confírma:
    ut, qui concéptum de Vírgine Deum verum et hóminem confitémur;
    per ejus salutíferæ resurrectiónis poténtiam, ad ætérnam mereámur perveníre lætítiam.
    Per eúndem Dominum\dots
}{
    Prosimy Cię, Panie, utwierdź w duszach naszych tajemnice prawdziwej wiary, abyśmy wyznając prawdziwe Bóstwo
    i człowieczeństwo Tego, który się począł z Dziewicy, zdołali dojść do wiekuistej radości radości mocą Jego
    zbawczego Zmartwychwstania. Przez tegoż Pana\dots
}

\oremus{Prefacja o Najświętszej Maryi Pannie}{
    Vere dignum et justum est, æquum et salutáre, nos tibi semper et ubíque grátias ágere: Dómine sancte,
    Pater omnípotens, ætérne Deus: Et te in veneratione beátæ Maríæ semper Vírginis collaudáre,
    benedícere et prædicáre. Quæ et Unigénitum tuum Sancti Spíritus obumbratióne concépit:
    et, virginitátis glória permanénte, lumen ætérnum mundo effúdit, Jesum Christum, Dóminum nostrum.
    Per quem majestátem tuam laudant Angeli, adórant Dominatiónes, tremunt Potestátes.
    Cæli cælorúmque Virtútes ac beáta Séraphim sócia exsultatióne concélebrant.
    Cum quibus et nostras voces ut admítti jubeas, deprecámur, súpplici confessióne dicéntes:
}{
    Zaprawdę godne to i sprawiedliwe, słuszne i zbawienne, abyśmy zawsze i wszędzie Tobie składali dziękczynienie,
    Panie, Ojcze święty, wszechmogący, wieczny Boże: I abyśmy Cię wielbili, błogosławili i wysławiali czcząc
    Najświętszą Maryję zawsze Dziewicę.
    Ona to poczęła Jednorodzonego Syna Twojego za sprawą Ducha Świętego i zachowując chwałę dziewictwa
    wydała światu światłość przedwieczną, Jezusa Chrystusa, Pana naszego.
    Przez Niego majestat Twój chwalą Aniołowie, uwielbiają Państwa, z lękiem czczą Potęgi.
    A wspólnie z nimi w radosnym uniesieniu sławią Niebiosa, Moce niebieskie i błogosławieni Serafini.
    Z nimi to, prosimy, dozwól i naszym głosom wołać w pokornym uwielbieniu:
}

\grecommentary[0.15cm]{Missa Cum iubilo (IX)}
\gregorioscore{Melody/ky--sanctus_ix--solesmes.1.gabc}

\medskip

\grecommentary[0.15cm]{Missa Cum iubilo (IX)}
\gregorioscore{Melody/ky--agnus_dei_ix--solesmes.1.gabc}

\oremus{Antyfona na Komunię (Iz 7,14)}{
    Ecce, Virgo concípiet et páriet fílium: et vocábitur nomen ejus Emmánuel.
}{
    Oto Panna pocznie i porodzi Syna i nazwą imię jego Emmanuel (Bóg z nami).
}


\oremus{Modlitwa po komunii}{
    Grátiam tuam, quǽsumus, Dómine, méntibus nostris infúnde: ut, qui, Angelo nuntiánte, Christi,
    Fílii tui, incarnatiónem cognóvimus; per passiónem ejus et crucem, ad resurrectiónis glóriam perducámur.
    Per eúndem Dominum\dots
}{
    Prosimy Cię, Panie, wlej w nasze serca Swoją łaskę, abyśmy poznawszy za zwiastowaniem anielskim
    wcielenie Chrystusa, Syna Twego, przez Jego mękę i krzyż zostali doprowadzeni do chwały zmartwychwstania.
    Przez tegoż Pana\dots
}

\newpage

\grecommentary[0.15cm]{Missa De Angelis (VIII)}
\gregorioscore{Melody/ky--ite_viii--solesmes.1.gabc}

\medskip

\proroctwomedskip{Antyfona}

\gregorioscore{Melody/an--alma_redemptoris--solesmes.1.gabc}

% \medskip

% \vfill

% \centerline{\pgfornament[width=1.5cm]{173}}

% \vfill

% \newpage

% \medskip

\vfill

\centerline{\Large \textbf{A. M. D. G.}}

\vfill

% \newpage

% \vfill

% \begin{tikzpicture}[color=black,
% every node/.style={inner sep=0pt}]
% \node[minimum width=12.8cm, minimum height=18.0cm] (vecbox) at (current page.center) {};
% \node[align=center, text width=10cm, inner sep=10pt] at (vecbox.center) {%
%     {\pgfornament[width=3cm]{173}}\\[1em]
% };
% \end{tikzpicture}


% \vfill


% \newpage

% \thispagestyle{empty}

% \vfill

% \begin{tikzpicture}[color=black,
% every node/.style={inner sep=0pt}]
% \node[minimum width=12.8cm, minimum height=18.0cm] (vecbox) at (current page.center) {};
% \node[align=center, text width=10cm, inner sep=10pt] at (vecbox.center) {%
%     {\Large \textbf{A. M. D. G.}}\\[1em]
% };
% \end{tikzpicture}

% \vfill

\end{document}
